\documentclass[sigplan,10pt,nonacm,review,anonymous]{acmart}
\settopmatter{printfolios=true,printacmref=false}
\setcopyright{none}
\renewcommand\footnotetextcopyrightpermission[1]{}
\acmDOI{}
\acmISBN{}

\usepackage{booktabs}
\usepackage{xspace}
\setlength{\emergencystretch}{1.5em}
% Generated by python -m agentipc.paper_numbers; do not edit.
\newcommand{\nMainTasks}{150\xspace}
\newcommand{\nMainRuns}{5400\xspace}
\newcommand{\nAccAOne}{72.3\xspace}
\newcommand{\nSentAOne}{13.3k\xspace}
\newcommand{\nCompAOne}{12.4k\xspace}
\newcommand{\nNotesAOne}{73\xspace}
\newcommand{\nDelivAOne}{99\xspace}
\newcommand{\nCallsAOne}{3.0\xspace}
\newcommand{\nAccAZc}{71.6\xspace}
\newcommand{\nSentAZc}{13.6k\xspace}
\newcommand{\nCompAZc}{4.8k\xspace}
\newcommand{\nNotesAZc}{72\xspace}
\newcommand{\nDelivAZc}{99\xspace}
\newcommand{\nCallsAZc}{3.0\xspace}
\newcommand{\nAccATwo}{62.1\xspace}
\newcommand{\nSentATwo}{5.6k\xspace}
\newcommand{\nCompATwo}{4.7k\xspace}
\newcommand{\nNotesATwo}{72\xspace}
\newcommand{\nDelivATwo}{99\xspace}
\newcommand{\nCallsATwo}{3.0\xspace}
\newcommand{\nAccB}{62.8\xspace}
\newcommand{\nSentB}{38.2k\xspace}
\newcommand{\nCompB}{9.8k\xspace}
\newcommand{\nNotesB}{59\xspace}
\newcommand{\nDelivB}{84\xspace}
\newcommand{\nCallsB}{17.0\xspace}
\newcommand{\nAccC}{66.3\xspace}
\newcommand{\nSentC}{37.4k\xspace}
\newcommand{\nCompC}{9.8k\xspace}
\newcommand{\nNotesC}{63\xspace}
\newcommand{\nDelivC}{93\xspace}
\newcommand{\nCallsC}{15.9\xspace}
\newcommand{\nAccD}{65.1\xspace}
\newcommand{\nSentD}{17.5k\xspace}
\newcommand{\nCompD}{7.2k\xspace}
\newcommand{\nNotesD}{50\xspace}
\newcommand{\nDelivD}{71\xspace}
\newcommand{\nCallsD}{7.8\xspace}
\newcommand{\nZcCut}{61\xspace}
\newcommand{\nSentRatioB}{2.9\xspace}
\newcommand{\nCompRatioB}{0.79\xspace}
\newcommand{\nSentRatioC}{2.8\xspace}
\newcommand{\nCompRatioC}{0.79\xspace}
\newcommand{\nSentRatioD}{1.3\xspace}
\newcommand{\nCompRatioD}{0.59\xspace}
\newcommand{\nSentRatioATwo}{0.4\xspace}
\newcommand{\nCompRatioATwo}{0.38\xspace}
\newcommand{\nGapAZcMin}{-1\xspace}
\newcommand{\nGapAZcMax}{3\xspace}
\newcommand{\nGapAZcSig}{0\xspace}
\newcommand{\nGapATwoMin}{3\xspace}
\newcommand{\nGapATwoMax}{15\xspace}
\newcommand{\nGapATwoSig}{5\xspace}
\newcommand{\nGapBMin}{6\xspace}
\newcommand{\nGapBMax}{14\xspace}
\newcommand{\nGapBSig}{6\xspace}
\newcommand{\nGapCMin}{-1\xspace}
\newcommand{\nGapCMax}{8\xspace}
\newcommand{\nGapCSig}{2\xspace}
\newcommand{\nGapDMin}{2\xspace}
\newcommand{\nGapDMax}{12\xspace}
\newcommand{\nGapDSig}{4\xspace}
\newcommand{\nZcRatioMin}{38\xspace}
\newcommand{\nZcRatioMax}{49\xspace}
\newcommand{\nFileSentMin}{1.8\xspace}
\newcommand{\nFileSentMax}{4.8\xspace}
\newcommand{\nPipeBrokeMax}{16\xspace}
\newcommand{\nPipeBrokeAOne}{2\xspace}
\newcommand{\nDelivBMin}{74\xspace}
\newcommand{\nDelivBMax}{86\xspace}
\newcommand{\nDelivCMin}{89\xspace}
\newcommand{\nDelivCMax}{97\xspace}
\newcommand{\nDelivDMin}{61\xspace}
\newcommand{\nDelivDMax}{78\xspace}
\newcommand{\nNotesCharsMsg}{1,559\xspace}
\newcommand{\nNotesCharsFile}{1,136\xspace}
\newcommand{\nNotesShorter}{27\xspace}
\newcommand{\nSamePrompt}{59\xspace}
\newcommand{\nDocReadsB}{83\xspace}
\newcommand{\nNoNotesB}{42\xspace}
\newcommand{\nDocReadsC}{38\xspace}
\newcommand{\nNoNotesC}{7\xspace}
\newcommand{\nHotpotUsable}{6,947\xspace}
\newcommand{\nWikiUsable}{11,281\xspace}
\newcommand{\nMaxLen}{32,768\xspace}


\newcommand{\figdir}{figs/out}
\newcommand{\ch}[1]{\textsf{#1}}

\begin{document}

\title[Pipes, Files, or Handles?]{Pipes, Files, or Handles? Revisiting IPC for LLM Agent Pipelines}

\author{Anonymous Author(s)}
\affiliation{\institution{Anonymous Institution}\country{}}

\begin{abstract}
LLM agents that work as a team pass results to each other in one of two ways: they copy the conversation
into the next agent's prompt, or they share files. This is the message-passing versus shared-memory
question of operating systems. We test it on a Researcher--Writer--Reviewer pipeline with Qwen3 models on
vLLM, changing only the channel and what it carries. We count cost as tokens sent, as prefill the engine
computes after prefix caching, and as serving time. With small payloads and Qwen3-8B, copying is the most
accurate channel. Zero-copy broadcast, which orders the prompt so the prefix cache reuses the shared
history, keeps copying's accuracy while computing \nZcRatioMin--\nZcRatioMax{}\% as much prefill, but
with each channel served alone, six tasks at a time, it saves no time. Shared files send
\nFileSentMin--\nFileSentMax{}$\times$ the tokens, because every access is a model call that resends the
conversation. Served alone, they take up to \nRFourTimeB{}$\times$ the time and generate up to
\nRFourGenB{}$\times$ the tokens. They also lose information on both sides. With the instruction's wording
matched, writing notes through a tool call costs facts and accuracy at \nRFourWSigWord of \nRFourSizesWord model sizes, and readers
without an index skip the notes. When three parallel Researchers split sixteen questions, plain files lose
\nGapBKSixteen points, and neither a larger KV cache nor twice the tool budget closes the gap. A 4B model
loses \nPSizeFourGapAZc points with zero-copy ordering.
\end{abstract}

\maketitle

\section{Introduction}

Multi-agent LLM systems split a task among roles: one agent gathers facts, another writes, a third
checks~\cite{autogen,metagpt,chatdev}. Most of them pass results by copying. AutoGen's group chat puts
every earlier message into every agent's prompt, and pipelines such as ChatDev forward each stage's
output to the next~\cite{autogen,chatdev}. A growing set of systems share files instead. Coding agents
read and write a workspace~\cite{sweagent,openhands}, and production research agents store subagent
results and pass references~\cite{anthropic-multiagent}. Prior work on agent communication cuts tokens by
pruning who talks to whom or by compressing what they say~\cite{agentprune,optima,pact}. It keeps the
channel fixed, so it cannot tell whether copying or sharing is the cheaper way to pass the same result.

Operating systems face the same choice. Message passing copies data into the
receiver. Shared memory avoids the copy but makes the receiver find and fetch what it needs, and the
classic finding is that neither wins everywhere~\cite{lauer-needham,liedtke-ipc}. For LLM agents the
costs look different. Each copied token is prefilled again. Each file read is a tool call, and because the
model is stateless, the call resends the whole conversation. The serving engine changes both costs,
because its prefix cache reuses the KV of a repeated prefix~\cite{vllm,sglang}.

This paper presents a controlled comparison of inter-agent channels. We keep the task, the model, the
roles, and their order fixed, and change only how the Researcher's and Writer's results reach the next
agent, over five channels (\S\ref{sec:channels}). Our key observation is that \emph{cost has to be counted
three ways}, because the channels trade one cost for another. Specifically, at three questions per task
shared files send \nSentRatioB{}$\times$ the tokens that copying the history does, which is what an API
bills when it does not discount cached input. Yet the engine computes only \nCompRatioB{}$\times$ as much
prefill. Serving time, measured with each channel served alone on one GPU at six tasks at a time, ranks
the channels as tokens sent does. Plain files take \nRFourTimeB{}$\times$ broadcast's time, and their
agents generate \nRFourGenB{}$\times$ the tokens. Copying the history with the shared part
first (zero-copy broadcast) cuts prefill by \nZcCut{}\% at copying's accuracy but takes
\nRFourTimeAZc{}$\times$ the time. Our second observation is that these results depend on the cache and the
model as well as on the channel. Prefill savings shrink when the agents' prompts no longer fit the KV cache,
and the file channels need a model that uses tools well.

The challenge is that changing the channel also changes what reaches the next agent. We therefore add
controls that change one without the other, for example by replaying the same notes into every channel.
We make three contributions:
\begin{itemize}
\item We build a testbed that varies the inter-agent channel, with such controls (\S\ref{sec:method}).
\item We measure five channels on two multi-hop QA benchmarks with \nMainTasks paired tasks per condition,
in tokens, prefill, and serving time, and find where each one loses information (\S\ref{sec:results}).
\item We vary payload size, pipeline length, tool budget, batched reads, and model size
(\S\ref{sec:sensitivity}), and derive rules for choosing a channel (\S\ref{sec:conclusion}).
\end{itemize}

\section{Five Channels}
\label{sec:channels}

\begin{figure*}[t]
\centering
\includegraphics[width=\textwidth]{\figdir/fig_mechanism.pdf}
\caption{How results reach the next agent. Top: message channels copy content into the next prompt (bars).
Zero-copy broadcast puts the shared part first, so the engine reuses it from the cache (hatched). Bottom:
file channels leave results in a shared workspace that later agents read through tool calls.}
\label{fig:channels}
\end{figure*}

The pipeline has three agents. The \emph{Researcher} reads the task's documents and writes down the facts
each question needs. The \emph{Writer} drafts a short answer and evidence per question. The \emph{Reviewer}
checks the draft and produces the final answers. Every Researcher receives the same user message, which
holds the questions and all documents. Its instructions differ in where its notes go (a reply, or files
written with a \texttt{write\_file} tool) and in what each channel asks of it, and
\S\ref{sec:results} separates the two for B. The channels differ in
what the Writer and the Reviewer receive (Figure~\ref{fig:channels}). Zero-copy broadcast is broadcast with
a different prompt order, so we count it as a variant of broadcast.

\paragraph{A1, broadcast.} Every earlier message is copied into the next prompt: the task with its documents,
the Researcher's notes, and, for the Reviewer, the draft. This is AutoGen's group chat~\cite{autogen}.

\paragraph{A1-zc, zero-copy broadcast.} The same messages as A1, but each prompt starts with the shared
history and ends with the role's instruction, which moves out of the system prompt. Consecutive agents then
share a prefix, and the engine reuses its KV instead of computing it again. Every token is still sent. The
new order also changes what agents write. At $k=3$ the Researcher's notes grow from \nNotesCharsAOne to
\nNotesCharsAZc characters.

\paragraph{A2, pipe.} Each agent receives only the previous agent's output: the Writer gets the notes, the
Reviewer gets the draft. ChatDev's chat chain works this way~\cite{chatdev}.

\paragraph{B, shared files.} The Writer and Reviewer receive only the questions. The documents and the
Researcher's notes are files in a shared workspace, read through \texttt{list\_files}, \texttt{read\_file},
and \texttt{write\_file} tools. Each tool call is one model call.

\paragraph{C, files with an index.} B plus \texttt{index.md}, which the workspace keeps up to date with each
file's author, size, and first line, as a file system keeps a directory. Agents are told to read it first,
and Researchers are told to head each file with its questions.

\paragraph{D, handles.} The Researcher and Writer write their files and also push a handoff of at most 120
words that gives each question's key finding and the path of its details. Later agents read files only for
what the handoff lacks.

\section{Methodology}
\label{sec:method}

\paragraph{Tasks.} We use HotpotQA in the distractor setting~\cite{hotpotqa} and
2WikiMultihopQA~\cite{2wiki}, the benchmarks prior agent-communication studies use~\cite{optima,pact}. Each
question comes with ten paragraphs and the sentences that support its answer. A task bundles $k$ questions
so the amount passed between agents grows with $k$. Its documents are the union of their paragraphs (about
30 at $k=3$). We drop yes/no questions and questions with fewer than two supporting sentences and draw \nMainTasks tasks per dataset and $k \in \{1,3,5\}$.
Every channel runs the same tasks, so comparisons are paired.

\paragraph{Model and serving.} Qwen3-8B~\cite{qwen3} with thinking disabled, greedy decoding, on vLLM
v0.33.0 with prefix caching, one RTX 3090 per server (\nMaxLen-token context). Tool calls use vLLM's
function-calling parser. Each server runs six tasks at a time, interleaved across channels, and its KV cache
holds \nBigKv tokens. Each run gets its own cache namespace (\texttt{cache\_salt}), so channels never share
cached KV. The timing, cache, and model-size runs use one RTX PRO 6000 per server
and a 100k-token KV cache unless stated.

\paragraph{Metrics.} \emph{Accuracy} is the share of questions whose final answer matches the gold answer
(exact match, F1 $\ge 0.8$, or the gold answer inside a short prediction). \emph{Tokens sent} is the sum of
prompt and completion tokens over all calls. It is what an API bills without caching and an upper bound
when the API discounts cached input. \emph{Tokens computed} counts prefill only, the prompt tokens the engine
did not serve from its cache. \emph{Serving time} is the wall time of a run that serves one channel alone
on one GPU, six tasks at a time. \emph{Fact delivery} asks which
supporting sentences a text states, scored with AlignScore~\cite{alignscore} at its 0.5
threshold. We score the Researcher's output and what the Writer received, and call the
share of the output's facts that the Writer received \emph{delivery}. In message channels the Writer
receives the copied notes, so delivery is complete by construction. In file channels it receives the text
of the files it read, source documents included. We attribute each wrong answer to one stage. If the
draft answer was right, we charge the Reviewer. Otherwise we charge the Researcher if its output misses a
supporting fact, delivery if the Writer did not receive one, and the Writer if it had them all. The split is
coarse, since \nRightMissAOne{}\% of broadcast's correct answers also come from notes that miss a supporting
fact.

\paragraph{Validating fact delivery.} On \nSynN pairs built from the task documents, with a supporting
sentence kept, removed, or with its key entity replaced, AlignScore separates the labels with an AUC of
\nSynAAuc, while Qwen3-8B used as a judge accepts \nSynQAcceptDropped{}\% of the removed sentences. Real
notes paraphrase, and on them AlignScore agrees only moderately with the judge ($\kappa$ = \nAsQKappa), so
we treat fact-level numbers as approximate. An exact match of each question's gold-chain entities, like
answer-string recall in retrieval~\cite{dpr}, agrees that notes in file channels keep fewer of the gold
entities. It
credits D's handoffs, which name the entities without stating the facts, with \nEntDDeliv{}\% delivery
where AlignScore finds about half.

\paragraph{Harness.} An agent may make up to 12 tool calls. If it stops without leaving its output file, it
gets one call that must write it, since frameworks take an agent's last message as its output. At $k=3$
this rescue is used in \nForcedMin--\nForcedMax{}\% of the runs of B's Writer and of the Reviewers of B
and C. Counting every rescued run as wrong would leave B, C, and D at \nStrictMin--\nStrictMax{}\%
accuracy. The Researcher writes but does not read, since its documents are already in its prompt. We report
95\% bootstrap intervals over tasks and paired differences against A1 unless stated. A range in parentheses after a difference is that interval, and counts of significant settings
are not corrected for multiple comparisons. Each task runs once per channel.

\section{Results}
\label{sec:results}

\begin{figure*}[t]
\centering
\includegraphics[width=\textwidth]{\figdir/fig_tradeoff.pdf}
\caption{Accuracy against cost per task, $k=3$, both datasets pooled (300 tasks per channel; bars are 95\%
intervals). Counted by tokens sent (a), copying the history is the accurate choice and files are the
expensive one. Counted by tokens computed (b), zero-copy broadcast keeps that accuracy at a pipe's prefill.
Labels in (b) give the share of prompt tokens that the engine served from its prefix cache.}
\label{fig:tradeoff}
\end{figure*}

\paragraph{Copying and zero-copy.} At $k=3$, with both datasets pooled, broadcast answers \nAccAOne{}\% of
the questions, more than any other channel (Figure~\ref{fig:tradeoff}). Zero-copy broadcast answers \nAccAZc{}\%, within \nGapAZcMax points of
broadcast in every dataset and $k$, and no difference is significant. Under exact match it is up to
\nZcEmGapMax points lower on 2Wiki. It computes \nCompAZc tokens per task to broadcast's \nCompAOne, close
to the pipe's \nCompATwo, because the Writer and the Reviewer reuse the documents the Researcher
already prefilled. Its tokens sent stay those of broadcast (\nSentAZc), so the gain is in prefill alone.

\paragraph{A pipe is cheap but loses at review.} The pipe sends the fewest tokens (\nSentATwo) but answers
\nGapATwoMin--\nGapATwoMax points fewer questions than broadcast, significantly in \nGapATwoSigWord of six
settings. The Reviewer sees only the draft and changes right answers to wrong
ones in up to \nPipeBrokeMax{}\% of the questions. In broadcast, where it also sees the notes and
documents, it does so in at most \nPipeBrokeAOne{}\% (Figure~\ref{fig:loss}b). On \nScopeATwoDN tasks, a pipe that
also passes the documents answers \nScopeATwoDVsATwoAbs points more than the pipe (\nScopeATwoDVsATwoLo to
\nScopeATwoDVsATwoHi). It still trails broadcast (\nScopeATwoDVsAOne points, \nScopeATwoDVsAOneLo to
\nScopeATwoDVsAOneHi) at \nScopeATwoDCompRatio{}$\times$ broadcast's prefill.

\paragraph{Files cost more tokens.} B and C send \nSentRatioB{}$\times$ and \nSentRatioC{}$\times$ the
tokens of broadcast at $k=3$, and \nFileSentMin--\nFileSentMax{}$\times$ across settings. Their pipelines
make \nCallsB and \nCallsC model calls to broadcast's three, and each call resends the conversation. The
cache absorbs most of it (Figure~\ref{fig:tradeoff}b), so B and C compute \nCompRatioB{}$\times$ the tokens of broadcast. D's handoffs
replace most reads, so it sends only \nSentRatioD{}$\times$ and computes \nCompRatioD{}$\times$.

\paragraph{Prefill savings buy no serving time.} We served each channel alone on one GPU with Qwen3-8B,
six tasks at a time. Broadcast takes \nRFourTimeAOne\,s for all \nRFourAloneN three-question tasks.
Zero-copy broadcast takes \nRFourTimeAZc{}$\times$ as long, and its agents generate \nRFourGenAZc{}$\times$
the completion tokens. B, C, and D take \nRFourTimeB{}$\times$, \nRFourTimeC{}$\times$, and
\nRFourTimeD{}$\times$ the time and generate \nRFourGenB{}$\times$, \nRFourGenC{}$\times$, and
\nRFourGenD{}$\times$ the tokens. They also make more model calls, in the same order, so for files we cannot separate generation from sequential
calls. Only the pipe is faster (\nRFourTimeATwo{}$\times$, with \nRFourGenATwo{}$\times$ the tokens). At
this load the channels rank by time as by tokens sent.

\begin{figure*}[t]
\centering
\includegraphics[width=\textwidth]{\figdir/fig_loss.pdf}
\caption{Where information is lost, $k=3$, both datasets pooled. (a) Shares of all supporting facts that
the Researcher's output states and that reached the Writer (AlignScore). The second bar is the first times
delivery, the share of the output's facts that reached the Writer.
(b) Wrong answers by stage at fault; an answer is
charged to the Researcher whenever its output misses any supporting fact, which favors that stage.}
\label{fig:loss}
\end{figure*}

\paragraph{Files lose information when written and when read.} Given the same questions and documents,
the Researcher writes \nNotesShorter{}\% shorter notes at $k=3$ when they go to files.
In B its notes state \nNotesB{}\% of the supporting facts, where broadcast's state \nNotesAOne{}\%
(Figure~\ref{fig:loss}a). On the reading side, the files B's Writer reads state only
\nDelivBMin--\nDelivBMax{}\% of the facts in the Researcher's notes. The file list does not say which file
holds the notes, so \nDocReadsB{}\% of B's Writer reads go to source documents, and in \nNoNotesB{}\% of the
tasks it never opens the notes. In D the Writer trusts the handoff, opens even fewer files, and receives
\nDelivDMin--\nDelivDMax{}\% of the notes' facts. The index names the notes. C's Writer skips them in
\nNoNotesC{}\% of the tasks, and delivery rises to \nDelivCMin--\nDelivCMax{}\%. In accuracy, C comes
within \nPooledGapC points of broadcast at $k=3$ pooled, and within \nGapCMax in any one
setting, while B trails it by \nGapBMin--\nGapBMax points in each of the six settings. C is significantly above B in \nCoverBSigWord of
six settings.

\paragraph{Giving files the documents.} Putting the documents in the prompts of B's later agents too
leaves B's accuracy unchanged on \nScopeBDN tasks (\nScopeBDVsB points, \nScopeBDVsBLo to \nScopeBDVsBHi)
at \nScopeBDSentRatio{}$\times$ its tokens. Missing sources therefore do not explain the
loss.

\paragraph{Writing through a tool call loses facts.} B's Researcher instruction differs from A1's in wording
as well as in the tool. Two controls on the \nRFourEightWMinusRN tasks of the model-size runs
(\S\ref{sec:sensitivity}) separate the two. In A1r the Researcher gets B's instruction with only its
destination sentences changed to a reply, and in A1w it gets B's instruction and writes with
\texttt{write\_file}. Both copy the notes forward as in broadcast.
A1r's notes state as many supporting facts as A1's at every size (\nRFourEightNotesR{}\% vs
\nRFourEightNotesAOne{}\% at 8B), and A1r is as accurate as A1 except at 4B, where it answers
\nRFourFourRMinusAOne points more. The wording thus costs no accuracy. A1w's notes state about as many facts as
B's (\nRFourEightNotesW{}\% vs \nRFourEightNotesB{}\% at 8B). At 4B, 8B, and 32B that is fewer than A1r's,
and A1w changes accuracy relative to A1r by \nRFourFourWMinusR points (\nRFourFourWMinusRLo to
\nRFourFourWMinusRHi), \nRFourEightWMinusR (\nRFourEightWMinusRLo to \nRFourEightWMinusRHi), and
\nRFourThirtyTwoWMinusR (\nRFourThirtyTwoWMinusRLo to \nRFourThirtyTwoWMinusRHi). At 14B the tool call
keeps the facts (\nRFourFourteenNotesW{}\% vs \nRFourFourteenNotesR{}\%) and A1w matches A1r
(\nRFourFourteenWMinusR, \nRFourFourteenWMinusRLo to \nRFourFourteenWMinusRHi). Writing through a tool
call, not the wording, costs accuracy whenever it drops facts from the notes.

\paragraph{Reading the same notes.} To measure reading alone we replayed broadcast's Researcher output,
unchanged, into every channel, as one file \texttt{notes.md} for B and C. Given the same notes, B and C
both answer \nRepBGap points fewer questions than broadcast on \nRepN tasks (\nRepBGapLo to \nRepBGapHi),
and the pipe \nRepATwoGap fewer. When the notes are a single, clearly named file, the index does not help.
C's Writer, whose index also lists thirty source documents, opens the notes less often than B's (\nRepOpenC{}\% vs \nRepOpenB{}\% of tasks) and receives fewer
of their facts (\nRepCDeliv{}\% vs \nRepBDeliv{}\%).

\section{Sensitivity}
\label{sec:sensitivity}

\begin{figure*}[t]
\centering
\includegraphics[width=\textwidth]{\figdir/fig_scaling.pdf}
\caption{Cost and accuracy relative to broadcast as the number of questions per task grows (both datasets
pooled). $k=16$ uses three parallel Researchers.}
\label{fig:scaling}
\end{figure*}

\paragraph{Larger payloads.} At eight questions per task each prompt holds about 80 paragraphs, so six
concurrent runs no longer fit the RTX 3090 servers' \nBigKv-token KV cache together. Zero-copy broadcast then
computes \nCompAZcKEight{}$\times$ broadcast's prefill, up from \nCompAZcKThree{}$\times$ at three
questions. B and C compute more than broadcast (\nCompBKEight{}$\times$ and \nCompCKEight{}$\times$),
because the prefixes they resend no longer stay cached between calls (Figure~\ref{fig:scaling}a). The pipe,
still cheap, now loses little (\nGapATwoKEight points). At sixteen questions, with three
parallel Researchers, zero-copy broadcast saves nothing (\nCompAZcKSixteen{}$\times$). The pipe
matches broadcast's accuracy within \nGapATwoKSixteen point at \nCompATwoKSixteen{}$\times$ its compute.
Plain files collapse. B answers \nGapBKSixteen points fewer questions than broadcast, C
\nGapCKSixteen fewer, and D \nGapDKSixteen (Figure~\ref{fig:scaling}b).

\paragraph{The cache does not explain the collapse.} A larger cache brings the savings back, fully at
eight questions and partly at sixteen.
At eight questions, changing only the KV cache on \nPKvThirtyFiveN tasks, zero-copy broadcast computes
\nPKvThirtyFiveCompAZc{}$\times$ broadcast's prefill at 35k tokens and \nPKvHundredCompAZc{}$\times$ at
100k, while accuracy moves by at
most \nPKvAccShiftMax points. At sixteen questions with a 100k-token cache (\nRFourKSixteenGapBN tasks),
zero-copy broadcast computes \nRFourKSixteenCompAZc{}$\times$ broadcast's prefill and B
\nRFourKSixteenCompB{}$\times$, against \nCompAZcKSixteen{}$\times$ and \nCompBKSixteen{}$\times$ on the same tasks with the RTX 3090
servers' \nBigKv-token cache. The accuracy gaps stay. B trails
broadcast by \nRFourKSixteenGapB points (\nRFourKSixteenGapBLo to \nRFourKSixteenGapBHi), D by
\nRFourKSixteenGapD (\nRFourKSixteenGapDLo to \nRFourKSixteenGapDHi), and C by \nRFourKSixteenGapC
(\nRFourKSixteenGapCLo to \nRFourKSixteenGapCHi), while the pipe stays within \nRFourKSixteenGapATwo
points (\nRFourKSixteenGapATwoLo to \nRFourKSixteenGapATwoHi).

\paragraph{A larger tool budget does not close the gap.} At $k=16$ B's Reviewer hits the 12-call budget in
\nLargeBReviewerLimit{}\% of runs, up from
\nKThreeBReviewerLimit{}\% at $k=3$. Yet in the separate 100k-cache runs, 24 calls change B's accuracy by only
\nRFourBTwentyFour points (\nRFourBTwentyFourLo to \nRFourBTwentyFourHi, \nRFourBTwentyFourN tasks), and
B's Reviewer still uses all 24 in \nRFourBTwentyFourRevLimit{}\% of runs (\nRFourBTwelveRevLimit{}\% with
12).

\paragraph{Parallel writers race on shared files.} What remains is specific to files at sixteen questions,
where three Researchers write in parallel and later agents search a larger workspace. Writing at the
same time, the Researchers pick the same file name. In \nClashB{}\% of B's tasks at least two write the same path, and the later write replaces
the earlier notes. B's notes keep \nLargeClashNotes{}\% of the supporting facts in the \nLargeClashN tasks
with a clash and \nLargeNoClashNotes{}\% in the \nLargeNoClashN without one. The race we see is not the
whole loss, since B answers \nLargeClashAcc{}\% of the questions with a clash and \nLargeNoClashAcc{}\%
without one. In C, whose Researchers head each file with its questions and see an index, no two wrote the
same path. With three Researchers on \nFanN
three-question tasks, B's writers clash in \nFanClashB{}\%, but B's accuracy change from one to three
Researchers, less broadcast's, is \nFanDiD{} points (\nFanDiDLo{} to \nFanDiDHi{}).

\paragraph{A fourth agent.} A Fact-checker before review adds \nChainAOneComp{}\% to broadcast's
prefill but \nChainAZcComp{}\% to zero-copy broadcast's, which it reuses. On
\nChainN tasks it raises the pipe from \nChainATwoAccRef{}\% to \nChainATwoAcc{}\% (\nChainATwoLo{} to
\nChainATwoHi{}) and lowers B from \nChainBAccRef{}\% to \nChainBAcc{}\% (\nChainBLo{} to
\nChainBHi{}).

\paragraph{Cheaper reads make the Writer skip the notes.} A \texttt{read\_files} tool that reads several
files in one call, the agent's \texttt{readv}, cuts B's calls from \nBbCallsB to \nBbCalls. But on \nBbN
paired tasks it lowers accuracy from \nBbAccB{}\% to \nBbAcc{}\% ($-$\nBbAccHi to $-$\nBbAccLo). Given a
cheap bulk read, the Writer reads the sources (\nBbDocReads{}\% of reads) and in \nBbNoNotes{}\% of
the tasks never opens the notes.

\paragraph{Model size.} With the cache at 100k tokens, we varied the Qwen3 model from 4B to 32B on
\nPSizeEightN three-question tasks. Files trail broadcast at every size, B by
\nPSizeFourteenGapB--\nPSizeFourGapB points and C by \nPSizeEightGapC--\nPSizeFourGapC, counting Qwen3-32B
with the path fix below. Qwen3-32B often writes to paths such as
\texttt{workspace/final.md}, so we reran it with a leading \texttt{workspace/} treated as the root. B, C,
and D then trail the first run's broadcast by \nRFourNormBGap (\nRFourNormBGapLo to \nRFourNormBGapHi), \nRFourNormCGap
(\nRFourNormCGapLo to \nRFourNormCGapHi), and \nRFourNormDGap points (\nRFourNormDGapLo to
\nRFourNormDGapHi), close to a rescore of the \nPRescueThirtyTwoBMoved B and \nPRescueThirtyTwoDMoved D runs
(of \nPSizeThirtyTwoN) whose answer the first harness missed. The pipe loses only with Qwen3-8B (\nPSizeEightGapATwo points,
\nPSizeEightGapATwoLo to \nPSizeEightGapATwoHi) and stays within \nPSizeFourGapATwo points at the other
sizes, so its loss at review is specific to that model. Zero-copy ordering costs only
the 4B model (\nPSizeFourGapAZc points, \nPSizeFourGapAZcLo to \nPSizeFourGapAZcHi). The loss is not the
longer notes. With broadcast's notes replayed into both channels, zero-copy still loses
\nRFourZcReplayFour points at 4B (\nRFourZcReplayFourLo to \nRFourZcReplayFourHi) and
\nRFourZcReplayEight at 8B (\nRFourZcReplayEightLo to \nRFourZcReplayEightHi). At 4B it comes from the
prompt order the Writer and Reviewer see.

\paragraph{Run-to-run noise.} Greedy decoding on a batching server is not deterministic. Two channels whose
Researcher receives the identical prompt produce identical output in \nSamePrompt{}\% of tasks. Rerunning A1
and B on \nAAN of the main tasks gives the same accuracy, with each difference's interval inside
$\pm$\nAAWorst points, and the same B--A1 gap (\nAAGapRepeat vs \nAAGapMain points).

\section{Related Work}

\paragraph{Agent communication.} AgentPrune, G-Designer, and AgentDropout prune the communication graph, and
Optima trains agents to send shorter messages~\cite{agentprune,gdesigner,agentdropout,optima}. PACT compares
what a message should contain~\cite{pact}. They keep the channel fixed, and we change it. Blackboard designs give agents a shared
state~\cite{lbmas,blackboard-ds}. MAST catalogs inter-agent failures~\cite{mast}, and our stage attribution
measures one of them, ignored input.

\paragraph{KV reuse across agents.} Prefix caching in vLLM and SGLang reuses an exact shared
prefix~\cite{vllm,sglang}, and Prompt Cache reuses attention states of prompt modules~\cite{promptcache}.
Parrot and KVFlow schedule multi-agent workflows so shared prefixes stay cached~\cite{parrot,kvflow}, and
KVCOMM and TokenDance reuse KV across agents whose prefixes differ by changing the
engine~\cite{kvcomm,tokendance}. Zero-copy broadcast applies the common practice of putting shared content first, so our contribution is
the measurement. TOAP
compares passing a reference, a summary, or a KV cache~\cite{toap}, not tool-mediated files.

\paragraph{IPC\@.} Lauer and Needham argued that message passing and shared memory are duals whose costs
depend on the system~\cite{lauer-needham}. L4 made messages cheap by cutting copies~\cite{liedtke-ipc}. In
agents, the prefix cache plays the role of shared pages, since every token is still sent, and the index
that of the directory.

\section{Conclusion}
\label{sec:conclusion}

The right channel depends on which cost the deployment pays, on payload size next to the KV cache, and
on the model. When the shared history fits the cache, copying it in an order the cache reuses kept Qwen3-8B
at broadcast's accuracy and computed \nZcCut{}\% less prefill at $k=3$. It saved no serving
time with each channel served alone at six tasks at a time, and it cost the 4B model accuracy. The pipe is
the cheapest channel by all three counts, but its savings come from dropping the documents, which at $k=3$ cost
Qwen3-8B's Reviewer. Long results are better written as plain replies that the system stores, since
writing them through a tool call lost facts and accuracy at \nRFourWSigWord of \nRFourSizesWord model sizes. Files need a model
that handles tools. An index helped Qwen3-8B find the notes among many files, most at sixteen questions,
but not Qwen3-14B and not when the notes were one named file. Plain files lost accuracy in every setting,
sent more tokens than copying, and took longer when served alone, though they computed
\nFileCompBMin--\nFileCompBMax{}$\times$ as much prefill at up to five questions. Our results come from one model family, linear pipelines, and multi-hop QA, not from
artifacts larger than a context window, such as code bases, where files may be the only option. We have not
tested why a tool call drops facts, why order matters at 4B, or what in file access at sixteen questions
causes the collapse.

\bibliographystyle{ACM-Reference-Format}
\bibliography{refs}

\end{document}
